% ES100 Project Proposal --- DESIGN-AGNOSTIC restructure (exploration copy).
% Sibling to proposal.tex (kept as reference). The change explored here:
% the offline/online split is a DESIGN decision --- in fact the axis along
% which the two design approaches differ --- so it is removed from the
% design-agnostic Description and relocated to the Design section, where the
% offline phase is characterized as part of Approach A.

% ======= LaTeX Setup ======= %
\documentclass[11pt]{article}
\usepackage[margin=1in]{geometry}
\usepackage{times}
\usepackage{indentfirst}
\usepackage{booktabs}
\usepackage{tabularx}
\newcolumntype{L}{>{\raggedright\arraybackslash}X}
\linespread{1.15}

\begin{document}

\section*{Project Proposal}

% ======= Problem Statement ======= %
\subsection*{Problem Statement}
The bandwidth of Bluetooth Low Energy (BLE) 6.0 is too low to support transmission
of raw neural extracellular electrophysiology data (hereby referred to as continuous
neural data) sampled at the industry standard from more than two recording channels
at once for offline spike sorting and consequent neuroscientific analysis [insert reference].

% ======= Need ======= %
\subsection*{Need Derived from Problem Statement}
Compress continuous neural data on the head stage of a rodent without domain-specific
loss of fidelity for wireless transmission.

% ======= Design Requirements and Technical Specifications ======= %
\section*{Design Requirements and Technical Specifications}
\subsection*{Description}

This section decomposes the stated need into the behaviors any valid solution must
exhibit (the functional requirements) and the measurable qualities those behaviors
must satisfy (the technical requirements). Both are stated independently of any
particular design; the design approaches that could realize them are considered
separately in the Design section.

Compressing continuous neural data without domain-specific loss of fidelity requires
isolating and preserving exactly the domain-relevant events---the spikes, their
identities, and their times---while discarding the remainder of the signal. Any
solution must therefore perform five operations: it must \emph{acquire} extracellular
voltage from multiple recording channels, \emph{detect} spike events within the
acquired signal, \emph{record} the time at which each detected spike occurs,
\emph{sort} each detected spike by assigning it a unit identity, and \emph{transmit}
each unit identity, paired with its timestamp, over a wireless link. These five
operations constitute the functional requirements of the system. Each imposes one or
more technical requirements---the measurable qualities its implementation must
satisfy---grouped by the operation that produces them in Table 1.

These technical requirements were first derived independently for each operation. One
requirement, however, resisted this per-operation decomposition: that no spike be
lost. A sample or detected spike dropped anywhere in the chain---whether during
acquisition, transmission to the sorter, or over the wireless link---compromises the
final result identically. It is therefore treated as a system-level requirement
spanning the full pipeline, summarized separately in Table 2, rather than attached to
any single functional requirement in Table 1.

\begin{table}[h]
\centering
\caption{Per-operation technical requirements of the system.}
\label{tab:perop}
{\footnotesize
\setlength{\tabcolsep}{4pt}
\renewcommand{\arraystretch}{1.25}
\begin{tabularx}{\textwidth}{@{}l L L L L@{}}
\toprule
Label & Requirement & Specification & Quantitative value & Measurement method \\
\midrule
T1  & $\geq 3$ channels amplified and digitized & channel count & $\geq 3$; target 32 & count channels in acquired stream \\
T2  & Sample fast enough for the spike band & per-channel sample rate & $\geq 12.5$~kS/s; design 25~kS/s & FFT of a known replayed signal \\
T3  & Digitization keeps up across channels & aggregate conversion rate & $\geq$ channels $\times$ per-channel rate & sustained conversion count \\
T4  & Resolve spikes above the noise floor & resolution / SNR & 16-bit; SNR [TBD] & noise and SNR on a known input \\
T5  & Link carries the digitized stream & ADC$\to$processor throughput & $\geq$ channels $\times$ rate $\times$ bits & dropped-sample count \\
T6  & Spikes separated from the LFP & high-pass corner & $\sim$300~Hz & filter frequency response \\
T7  & Detect spikes reliably & sensitivity / false-positive rate & TPR [TBD]; FPR [TBD] & compare to Kilosort 4.0 \\
T8  & A timestamp is recorded per spike & timestamp presence & one per detected spike & inspect output stream \\
T9  & Each spike assigned a unit identity & unit identity & one of $K$ units & inspect output \\
T10 & Sort better than chance & accuracy vs.\ Kilosort & $> 1/K$ (chance); [TBD] & confusion matrix vs.\ Kilosort \\
T11 & Packet carries identity and timestamp & packet content & 5~B (1~B ID $+$ 4~B ts) & inspect received packets \\
\bottomrule
\end{tabularx}
}
\end{table}

% [NEXT STEP: Table 2 --- system-level (cross-cutting) requirements:
%  no-loss/keep-up, and timestamp fidelity (stable clock, resolution,
%  stamped-at-detection).]

% ======= Design ======= %
\section*{Consideration of Multiple Design Approaches}

The functional and technical requirements above admit more than one design. The
approaches considered here differ principally in \emph{how} and \emph{when} spikes are
sorted---the operation that assigns each detected spike a unit identity. This single
choice determines whether the system requires a separate offline phase at all.

\subsection*{Approach A: Offline-Supervised Sorting (Two-Phase)}
This approach operates in two phases under different constraints. An offline phase,
run once on a workstation, sorts a reference recording via unsupervised clustering and
trains a classifier on the resulting labeled units. An online phase, running on the
SoC, performs supervised inference against this pre-trained model. The offline phase
is not subject to the power, latency, or memory constraints of the device; the online
phase is. Its requirements therefore include the offline preparation as a
precondition: a labeled unit library and trained classifier must exist before the
device operates. [Systematic statement of the offline-phase operations and their
requirements belongs here.]

\subsection*{Approach B: Online-Unsupervised Sorting (Single-Phase)}
This approach performs all sorting on the device as data arrive, using an online
clustering algorithm (e.g., Osort). No offline phase and no reference recording are
required; the system is self-contained and needs no per-subject calibration. Because
clustering must proceed without the full data set available in advance, sorted-unit
stability and accuracy are expected to be lower than the supervised approach---the
central trade-off between the two designs.

% [NEXT STEP: comparison of the two approaches against the critical technical
%  requirements (accuracy, memory, calibration burden, self-containment), and
%  the subsystem interactions each implies.]

\end{document}
